\section{Conclusions}
\label{sec:conclusions}

We have written a fresh Physical Review B draft directly from the archived Fe--Mo machine-learning repository and shown that the underlying workflow supports a data-efficient description of complex TCP intermetallics. The curated notebooks report 292 usable structures after curation, 261 non-bcc/fcc/hcp structures entering the stratified learning analysis, and a 70-structure complex-phase validation set spanning $R$, $P$, $M$, and $\delta$. Within that framework, domain-informed descriptors reduce the internal test errors from more than \SI{150}{\milli\electronvolt\per\atom} for atomic baselines to about \SI{26}{\milli\electronvolt\per\atom} for ACE/KRR and \SI{28}{\milli\electronvolt\per\atom} for the bond-specific BOP/KRR model, while coordination-number-resolved averaging lowers the BOP and SOAP errors by roughly one quarter to two fifths. The downstream validation and thermodynamic notebooks then show that these learned energies remain useful outside the immediate training manifold by supporting complex-phase parity analysis and finite-temperature $R$-phase occupancy predictions at \SI{1700}{\kelvin}. The physically motivated combination of chemistry-aware initialization, crystallography-aware aggregation, and bonding-aware descriptors therefore provides an efficient route to complex intermetallic energetics from a modest DFT dataset. The natural next steps are to incorporate magnetic corrections explicitly, to recover and tabulate the full external-validation metrics alongside the notebook figures, and to extend the same domain-guided workflow to ternary transition-metal systems and CALPHAD-coupled thermodynamic assessments.
